\section{Discussion and limitations}
\label{sec:limits}

\paragraph{What the evidence supports.}
On HELM MMLU the aggregation rule changes the price of first place by an order of magnitude, the ranking of the rules by price depends on whether the price is counted in points or in subjects, and the realism of the cap decides whether first place is reachable at all. The cheap rules (plurality, the minimum, the median) are the unstable ones under evaluation noise and task resampling. On the Open LLM Leaderboard the minimum is again cheap, but the savings of the median and plurality shrink, Borda becomes cheaper, the median is again unstable (relative width of the noise interval 1.83), and the order of the rules differs ($\tau=0.63$ in points, 0.46 in subjects), so we find no single cheapest rule and no order of rules that transfers between leaderboards, nor a law in which the cost scales with the number of sub-datasets (Section~\ref{sec:kn}).

\paragraph{Assumptions that carry the results.}
(i) \emph{Caps.} What a developer can achieve on a sub-dataset is not observable. Cap 100 and the best existing score show the price when scores are free to move, ``base $+h$'' how quickly first place becomes unreachable when they are not; no cap is claimed to be the truth, and a result is called robust only if it holds across caps. (ii) \emph{Single actor, static field.} Rivals do not react; several developers doing the same would raise the price. (iii) \emph{No transfer, no cost model.} A gain on one subject does not move another, and the cost of obtaining a point is not modelled: $L_1$ and $K$ are costs in score space, not in data, compute or money. (iv) \emph{Margin.} First place means strictly above every rival by $10^{-3}$ points, which is far below the evaluation noise of a subject with a few hundred items (Section~\ref{sec:robust}); the cost is a statement about the observed table.

\paragraph{Where the conclusions fail or are weak.}
(1) The order of the rules in $K$ is not stable across caps ($\tau$ between 0.19 and 0.87). (2) Medians under tight caps are over different feasible subsets (Section~\ref{sec:feas}); only the paired ratios compare rules, and they use only the targets feasible under both. (3) Headline medians are over all targets and are dominated by models far from first place; the cost for a contender (Section~\ref{sec:rank}) is much smaller. (4) The Condorcet-winner cost is an upper bound for other Condorcet-consistent rules; we computed only Copeland among them. (5) The weighted means use five random weightings, not any published one. (6) HELM subjects have 100--1534 items each, so single-subject scores are noisy; we studied noise and task resampling for ten targets at one cap (base $+10$) and do not extrapolate to other caps. (7) Only two leaderboards are studied; on the Open LLM Leaderboard the rank rules are sampled and medians over different target sets are not comparable, and the median target is a model far from first place.

\paragraph{Declared incomplete work.}
The whole analysis took about 15 hours of job time on two CPUs each (\texttt{code/REPRODUCE.md}), more than the 4 CPU-hour guidance of the task, although we scaled it down: on the Open LLM Leaderboard Borda, plurality and the Condorcet winner are computed for a rank-stratified sample of 57 targets, and Copeland (30 targets at cap 100, 116 at base $+10$) and the trimmed mean (cost in points only) for a top-500 slice only; 2 of the 30 and 2 of the 116 Copeland programmes reached the 120\,s limit and are reported as failed, not as infeasible. The raw-score variants use four caps. We did not compute Kemeny, Schulze or ranked-pairs rules, train any model, or check the single-actor model against real fine-tuning beyond the empirical per-task gains used as caps.
